% =============================================================================
%  3.1 Система реактивности
% =============================================================================
\section{Система реактивности}

\term{Реактивность}{Reactivity} в Vue~3 реализована через механизм
\inl{Proxy} (ES2015), который перехватывает операции чтения и записи
свойств объекта.  Это позволяет автоматически отслеживать зависимости
и обновлять только затронутые части интерфейса.

В Vue~2 использовался \inl{Object.defineProperty}, что накладывало
ограничения: невозможно было отследить добавление и удаление свойств.
Proxy снимает эти ограничения.

\subsection{Модель зависимостей}

Система реактивности строит ориентированный граф зависимостей.
Если свойство $p$ читается в эффекте $e$, создаётся ребро
$p \to e$.  При изменении $p$ все связанные эффекты помечаются
как «грязные».

\subsubsection{Сложность обновления}

При изменении свойства, от которого зависят $k$ эффектов:
\begin{equation}
  C_{\text{update}} = O(k), \quad k = |\{e : p \to e\}|
\end{equation}

\begin{code}{javascript}
import { reactive, watchEffect } from 'vue';

const state = reactive({
  count: 0,
  label: 'Счётчик',
});

watchEffect(() => {
  // автоматически отслеживает state.count и state.label
  console.log(`${state.label}: ${state.count}`);
});

state.count++;  // → «Счётчик: 1»
\end{code}

\note{Функция \inl{reactive()} работает только с объектами.
Для примитивных значений используйте \inl{ref()}, который
оборачивает значение в объект со свойством \inl{.value}.}

\important{Деструктуризация реактивного объекта разрывает реактивную
связь.  Используйте \inl{toRefs()} для сохранения реактивности
при деструктуризации: \inl{const \{ count \} = toRefs(state)}.}
